\section{Introduction}
One key enabler for the application of \ac{AI} and \ac{ML} is the establishment of standards and regulations for practitioners to develop safe and trustworthy \ac{AI} systems. Despite several efforts by different research institutions and governmental organizations, there are  currently no established guidelines. The most influential attempt stems from the European Union with its AI Act \citep{AiAct}.  
In this paper we revisit  the state of the art in safeguarding and certifying AI applications with a focus on the the European Union. Despite this, our summary is intended to be as broad as possible. The focus on European legislation stems from the influence on regulation and standardization it has.\footnote{\href{https://time.com/6278144/european-union-artificial-intelligence-regulation/}{European Union Set to Be Trailblazer in Global Rush to Regulate Artificial Intelligence (TIME)}} 
The result is a structured guideline for reporting AI applications along four major development steps.
The objective of our guideline is to provide comprehensive support during the development process. It encompasses a synthesis of requirements derived from both non-technical users and prominent certification and standardization bodies. Our approach involves the inclusion of robust references to relevant support materials and toolboxes, emphasizing the transparency and traceability of each requirement's source. For this purpose, we conduct interviews with standardization experts, developers and consumers of the final AI product and incorporate their feedback into our proposal. Furthermore, our framework serves as a solid foundation for prospective system certification. By adhering to our guidance, developers can ensure they are adequately equipped to meet forthcoming legal obligations within their domain.
The key contributions of our work is fourfold: 
\begin{itemize}
    \item Our guideline is meant to be used during the development process. 
    \item We combine requirements from non-technical users as well as certification and standardization bodies.
    \item We include references to support material and toolboxes and link the source of the requirement whenever possible.
    \item We lay the foundations for a future certification of the system. With our framework, developers should be well-positioned when legal requirements  come into force.
\end{itemize}

\subsection{Safeguarding AI, Law, Certification, Standardization, and Audits}
Before going into the specifics of AI regulation, we first give a short summary of relevant terms in EU regulations.
The term \emph{audit} describes the process of testing whether a product, service, or process meets certain requirements.
The result of such a process can be a \emph{certification}, the assurance by a third-party auditor that given requirements are fulfilled.
These requirements are often listed in \emph{harmonized standards}, which are created by organisations like CEN, CENELEC, or ETSI, following a request from the European Commission.
These harmonized standards provide the technical details necessary for companies to be able to ensure compliance with regulation and for auditors to have clear specifications to check against.
The process of creating such harmonized standards for the EU AI Act is currently running in parallel to the legislation procedure.
As a baseline for the standards harmonized with the EU AI Act, standards provided by ISO or other standardization bodies can be considered.
Conformity assessment bodies can then utilize these standards to showcase the compliance of products, services, or processes with applicable EU legislation \citep{EU-harmonised-standard}.
These audits can also be carried out on a voluntary basis for products that do not necessarily require them, for instance for marketing purposes.
In certain product categories like medical devices, third-party conformity assessments are mandatory to ensure compliance with technical regulations before entering specific markets \citep{EU-conformity-assessment}.
Apart from such certificates, so called \emph{guidelines} have no legally binding effect.

Finally, developing safe and responsible AI also involves governance structures within the company to manage data development activities and risks. However, this is beyond the scope of this paper, in which we focus on reporting around the development process. 

\subsubsection{Challenges in AI Certification}
Despite being a  topic of active research for several year, there is a lack of experience and methods for securing and certifying \acsu{ML} applications.
Apart from the fast progress in  \acsu{ML} around research and development, the challenges are rooted in the fundamental programming differences between classic software and \acsu{ML}. 
Classic software components are usually a combination of manually defined functions, the logic within these functions being explicitly specified by a human. In contrast, the core of \acsu{ML} is to learn from data while the user only specifies the model framework. Often, the result is a highly complex black-box that returns results without explaining the underlying decision-making process. 
The main challenge of certifying such a system is validating a black-box model representing a highly complex function.  Furthermore, an \acsu{ML}  algorithm  always entails risks related to data, e.g., the risk of a selection bias \citep {Tidjon2022_ChallengesAI}. %In addition, it’s important to note that the training data may not be representative of all possible scenarios and conditions that the model may encounter going live.  
Hence, within the certification process, not only does the proper functionality of the model need to be verified but also the representability of the data for training and testing. This is a challenging task, because the reason for developing \acsu{ML} often is that the underlying problem and cause-effect relations are too complex to be understood, formulated and solved analytically and manually.

\subsection{The State of \ac{AI} Certification}
There are several attempts by different organizations to standardize and secure the application of \ac{AI}. Apart from several ISO standards which are currently under development (e.g. \cite{ISO42001} - Artificial intelligence - Management system)  the European \ac{AI} Act \citep{AiAct} and the AI Assessment Catalog by Fraunhofer IAIS \citep{poretschkin2021_IAIS_Leitfaden} emerged during our discussions with organizations working on AI standardization and as such form a baseline for this work. They are described in more detail in this section.

%\pagebreak
\subsection{Non-academic Activities}
\subsubsection{The European \ac{AI} Act}
The most impactful work comes from the \ac{EU} with the proposal of the \ac{AI} Act~\citep{AiAct}. The first draft was proposed in April 2021 by the European Commission and two comments were recommended in December 2022 from the \cite{AiAct_Ratskompromiss} and in April 2023 by the \cite{AiAct_Parlamentsvorschlag}. 
The initial draft already created a heated debate leading to fears that it might over-regulate research and development of \ac{AI} in the \ac{EU} and that the definition of \ac{AI} was too broad. However, this debate is beyond the scope of this paper. We are interested in the requirements posed on those systems. Although some are still subject to change, many requirements have already emerged. 

In general, the \ac{AI} act distinguishes three risk classes for \ac{AI}: forbidden applications, high-risk applications and uncritical applications. Furthermore, there might be some obligations of transparency even for uncritical applications. For example, that end-users must be informed that they are interacting with an \ac{AI} system. There are no mandatory obligations for uncritical applications---although developers of such \ac{AI} systems shall still be motivated to follow the proposed principles. 

Whether an application falls into the high-risk category is risk-dependent, and a few high-risk applications are named in the \ac{AI} Act's appendix. Some of them include: biometric identification, operation of critical infrastructure, education, employment \& workers management, essential private services \& public services, law  enforcement, migration \& asylum, and administration of democratic processes. The requirements for such systems are listed in Chapter~2 and include 
\begin{itemize}
    \item a risk management system (Article 9)%
    \item  data and data governance (Article 10)%
    \item  technical documentation (Article 11)%
    \item record-keeping (Article 12)%
    \item  transparency and provision of information to users (Article 13)%
    \item  human oversight (Article 14)%
    \item  accuracy, robustness and cybersecurity (Article 15). 
\end{itemize}

A further distinction is based on whether an AI system is considered as \emph{general purpose AI} (\emph{GPAI}) or \emph{foundation model}. Although the term GPAI lacks a proper and widely accepted definition, GPAI should refer to systems such as image and speech recognition that  can by applied to a broad variety of tasks. Foundation models, on the other hand, will most likely be distinguished by the training data and the generality of the output. This should specifically target systems such as large language models like ChatGPT and pose stricter regulation on them while GPAI systems must only comply with the regulations in the Articles 9--15 if they are expected to be used in high-risk scenarios. The definitions for GPAI and foundation models (and others) are provided in Article 3. Finally, Chapter 3 lists further obligations for providers of high-risk AI systems and will most likely contain the requirements for foundation models once the AI Act is final.
%Some additional requirements are also listed in Chapter~3. A detailed discussion of these requirements is beyond the scope of this paper. 
At the moment, it is highly likely that a self-assessment will be possible for most applications. This means that companies can audit their applications by themselves and then apply the \href{https://en.wikipedia.org/wiki/CE_marking}{CE-mark} to their product. Third party assessments will only be mandatory for certain highly critical domains such as medical technology. 
Overall, the requirements are on a very abstract level and still subject to change. In particular, the act does not contain a specific starting points for practical implementation in companies because the baseline standards that will provide the technical details are still being developed.  

\subsubsection{\ac{AI} Assessment Catalog}
The \emph{\ac{AI} Assessment Catalog} is the most comprehensive and technical guideline with over 100 pages of content. It gives authors and practitioners a structured tool to develop and asses trustworthy \ac{AI} systems \citep{poretschkin2021_IAIS_Leitfaden}. The catalog considers the whole \ac{AI} application, i.e., the data and the AI components together with surrounding non-AI software components, and identifies six risk dimensions that need to be considered: fairness, autonomy and control, transparency, reliability, safety, and privacy. For each of these six dimensions there are guidelines for risk analysis and appropriate actions. The actions for each risk belong to one of the four groups: data, \ac{AI} components, embedding, or operations. Thus, if developers want to check whether they have addressed all necessary actions for data they have to consult all six risk dimensions. To the best of our knowledge, the AI Assessment Catalog is the most comprehensive guideline for trustworthy AI. However, due to its size the catalog can be cumbersome for practitioners to use. Finally, although being developed by a German organization, the guideline is not specific to Germany nor the \ac{EU} but combines elements from different areas of research and development.  

\subsection{Research Activities}
In addition to government and industry activities, there is research work on model and data reporting with cards that finds application especially among larger corporations.\footnote{\href{https://ai.googleblog.com/2022/11/the-data-cards-playbook-toolkit-for.html}{The Data Cards Playbook: A Toolkit for Transparency in Dataset Documentation}} \footnote{\href{https://ai.googleblog.com/2020/07/introducing-model-card-toolkit-for.html}{Introducing the Model Card Toolkit for Easier Model Transparency Reporting}} Our own work builds on documentation along these cards. 

\subsubsection{Model Cards}
\cite{Mitchell2019_ModelCards}  propose model cards as a standardized way to document \ac{ML} models, their performance metrics and characteristics. The work focuses primarily on ethics and fairness, citing that AI models in recent history have often erred more on social groups that have historically been marginalized in the US. All involved stakeholders benefit from a standardized use of model cards. The authors list nine sections for model cards:  model details,  intended use, factors, metrics,  evaluation data,  training data,  quantitative analyses,  ethical considerations, and caveats and recommendations.
In our proposal we divide model and data cards and therefore exclude the section on data. 

\subsubsection{Data Cards}
Data cards are similar to model cards \citep{Pushkarna2022_DataCards}. Their goal is to report several standard facts about a dataset in a simple and standardized form. The authors list several requirements for their cards. In particular, they should be standardized in order to be comparable, should be created at the same time as the dataset, written in a simple language such that users with a non-technical background can understand them, and, finally, report \emph{known unknowns} such as possible shortcomings or uncertainties of the dataset. Apart from that, the authors name 31 content themes to describe a dataset. They are listed in Appendix \ref{app:datasheets_content_themes}. 

Another related work is data sheets for datasets \citep{Gebru2021_DataSheets}. The authors propose a similar set of questions for dataset creators during the collection phase that can be grouped into the following categories: (i) \emph{motivation} for creating the dataset, (ii) \emph{composition} of the data, (iii) \emph{collection process}, (iv)  \emph{preprocessing/cleaning/labeling} of the data, (v) \emph{Uses} (describing the contexts in which the dataset was already used), (vi) the \emph{distribution} of the data, and (vii) the future \emph{maintenance}. The intended audience are both the dataset creators and the consumers of the final dataset.

\subsubsection{Foundation Models \& AI Act}
Recently, \cite{bommasani2023eu-ai-act} assessed popular foundation models like GPT-4, LLaMA and Luminous for their compliance with the AI Act. They derived and summarized 12 requirements grouped into data, compute, model and deployment from the AI Act. Our own approach is similar to their work but is not focused on foundation models. However, many of the requirements for foundation models also apply to high-risk AI systems. A summary of the requirements can be found in Table \ref{tab:stanford_llm_audit}.

\renewcommand\arraystretch{1.3}
\begin{table}
    \centering
    \begin{longtable}{p{0.15\linewidth}|p{0.6\linewidth}p{0.15\linewidth}}  
        \hline
         \textbf{Category}  & \textbf{Requirement}                                              & \textbf{Source}  \\
         \hline \hline
         \multirow{3}{=}{Data}
                            & Describe data sources                                             &  Annex VIII \\
                            & Apply data governance                                             &  Art. 28 b)\\
                            & Disclose copyrighted data                                         &  Art 28 b) \\
         \hline
         \multirow{2}{=}{Compute}
                            & Disclose compute (computer power, training time)                  &  Annex VIII \\
                            & Measure energy consumption                                        &  Art 28 b) \\
         \hline 
         \multirow{6}{=}{Model}
                            & Describe capabilities and limitations of model                    &  Annex VIII \\
                            & Describe foreseeable risks and associated mitigations             &  Annex VIII, Art. 28 b) \\
                            & Benchmark model on standard benchmarks                            &  Annex VIII, Art. 28 b)  \\
                            & Report results of internal and external tests                     &  Annex VIII, Art. 28 b)  \\
         \hline 
         \multirow{4}{=}{Deployment}
                            & Disclose content generated from the model                         &  Recital 60 g \\
                            & Disclose EU member states where model is on market                &  Annex VIII \\
                            & Provide documentation for downstream compliance with AI Act       &  Art 28 b) \\
         \hline 
    \end{longtable}
    \caption{Requirements for foundation models from \cite{bommasani2023eu-ai-act}.}
    \label{tab:stanford_llm_audit}
\end{table}

\subsubsection{Other Activities}
\label{sec:other_activities}
Besides the European \ac{AI} Act there are several other regulatory attempts. % regarding \ac{AI} standardization. 
The EU has been actively involved in standardization activities related to \ac{AI} and plans to implement liability rules on \ac{AI} to renew the existing rules on product liability. The new rules will hold manufacturers and importers of AI-powered products accountable for any defects or malfunctions that may cause harm to consumers \citep{EULiability}. Additionally, there are a number of ISO standards either in active development or already in use that target \ac{AI} systems in particular. The ISO/IEC JTC 1/SC 42 is the technical committee within the \ac{ISO} that is responsible for developing standards for \ac{AI}. Currently, there are 57 standards planned, some of which are already published. An overview can be found on the web page of the committee.\footnote{\href{https://www.iso.org/committee/6794475/x/catalogue/p/1/u/1/w/0/d/0}{Standards by ISO/IEC JTC 1/SC 42 : Artificial intelligence}}
Among the most important ones are: the \cite{ISO5259} series on data quality, \cite{ISO5338} on AI system life cycle processes, \cite{ISO8183} on a data life cycle framework, \cite{ISO23894} on risk management for AI systems and, \cite{ISO42001} for AI management systems. The last one is not addressed in this work because it deals mostly with governance, which we do not cover here. These norms emerged from discussions with certification bodies that we held in the course of our project.

Another recent and closely related work introduces a similar guideline to ours \citep{Wittenbrink2022_LeitfadenQualitatsmanagementKI}. It breaks down the AI development process into four steps and provides requirements as well as further instructions such as relevant norms for each requirement. The main difference to our work is, that the development process is broken down into characterisation, design, development, and deployment, whereas we follow an approach oriented on existing work on model and data cards. 

On the governmental side, China and the US are among the most important other players in regulating AI. The US approach significantly differs from the European one and is (currently) based on voluntary sector-specific measures. The most important non sector-specific work is the AI Risk Management Framework developed by the National Institute of Standards and Technology \emph{(NIST)} \citep{NISTAIRMF_2023}. It defines seven key characteristics for trustworthy AI that include (i) reliability, (ii) safety, (iii) security and resiliency, (iv) accountability and transparency,  (v) explainability and interpretability, (vi) privacy, and, (vii)  fairness. The companion \href{https://airc.nist.gov/AI_RMF_Knowledge_Base/Playbook}{NIST AI RMF Playbook} names four pillars to mitigate these risks: (i) map: risks are identified, (ii) measure: the identified risks are analyzed and tracked, (iii) manage: risks are prioritized and acted upon and, finally, (iv) govern: a culture of risk management and governance is present. The key characteristics of trustworthy AI are well aligned to those defined in the AI Act.\footnote{\href{https://www.nist.gov/system/files/documents/2023/01/26/crosswalk_AI_RMF_1_0_OECD_EO_AIA_BoR.pdf}{An illustration of how NIST AI RMF trustworthiness characteristics
relate to the OECD Recommendation on AI, Proposed EU AI Act, Executive Order 13960, and Blueprint for an AI Bill of Rights}} Furthermore, NIST is working to 
 align the framework with international standards that are also relevant for the AI Act.\footnote{\url{https://www.nist.gov/artificial-intelligence/technical-ai-standards}}

China, on the other side, follows a vertical approach\footnote{\href{https://carnegieendowment.org/2023/02/14/lessons-from-world-s-two-experiments-in-ai-governance-pub-89035}{Lessons From the World’s Two Experiments in AI Governance (carnegieendowment.org/)}} that targets each AI applications with a specific regulation and aims to be a leader in AI standardization \citep{Huw_2021_ChinaAIRegulation}. For example, one regulation specifically targets services that generate text, images or videos. Overall, although the high level pillars are expected to be similar to the ones in western countries and include that AI should benefit human welfare, China's regulation is expected to be focused more around group relations than individual rights \citep{Huw_2021_ChinaAIRegulation}. However, this is beyond the scope of our paper and Chinese regulation is not addressed in our framework.

\subsection{Contributions and Methods}
We merge contributions of all of the previous works into a single framework for reporting AI applications. As mentioned, this work focuses on the documentation of the AI application and does not consider governance  within a company. 
Our guideline follows the development process of an AI system with the goal to support developers during various model development steps in creating trustworthy AI systems from the beginning and positioning them well when legal requirements come into force. Apart from this, we cover the same risk dimensions as previous approaches and reference the AI Act and other guidelines whenever possible. The AI development cycle is broken down into four steps: (i) the use-case definition, (ii) data collection, (iii) model development, and finally,  (iv) model operation. We use four cards as method of reporting in order to reduce the number of resulting documents so  the hurdle to access them is as low as possible while the different topics remain separated. In the coarse of the project we furthermore evaluated the approach on three use-cases and integrated both feedback from developers and other affected individuals. Finally, we also held interviews with several certification bodies to get an understanding about the most important existing documents and guidelines for our work.

\subsubsection{Motivation for Cards}
Apart from discussions with certification experts preparing for the certification of AI systems, we also polled employees being affected by the introduction 
of AI systems. Our goal was to assure that our proposal would both satisfy the requirements of possible auditors and builds trust by the individuals being affected by the AI solution. 
The affected individuals mainly wished to be included in the development process in order to convince themselves about the functionally of the AI system. Furthermore, the system should be transparently visualised and the output and reasons for the decision should be comprehensibly explained. %Regarding certification, the partners confirmed that a certification would increase their trust as long as it is performed by a trustworthy institution. 
The reason why we choose the format of cards that follow the development process instead of following risk dimensions (in contrast to, for instance, \cite{poretschkin2021_IAIS_Leitfaden}) is twofold: firstly, for developers it is easier to follow the development process and for domain experts and other employees the resulting number of documents is smaller and as a result easier to access.
Secondly, although AI development is often an agile process, the cards can still be utilized by updating them when the dataset changes or the model is updated. They also integrate well into a classical v-model development, which is subject to future work. 

